This project set out to develop a low-cost multimodal drone detection 
system using exclusively off-the-shelf and 3D-printed components, 
combining acoustic and visual sensing to achieve reliable detection 
without specialist hardware. The results demonstrate that this goal 
was largely achieved within the project constraints.

The acoustic subsystem successfully progressed through three 
microcontroller iterations before arriving at the STM32F413 with 
DFSDM-based PDM-to-PCM conversion as the final working 
configuration. The GCC-PHAT direction of arrival algorithm 
demonstrated accurate azimuth estimation of a live drone in 
laboratory conditions, with a full pipeline iteration time of 
approximately 1.285 seconds including sound classification. 
The CNN sound classifier achieved 99\% accuracy on microphone 
array data and performed reliably during live testing, correctly 
distinguishing drone sound from background noise in all but one 
outlier iteration out of thirty. The decision-tree accumulative 
classifier, while performing well on public dataset evaluation, 
failed to generalise to the microphone array recordings, indicating 
a domain mismatch between training data and real deployment 
conditions that warrants further investigation.

The visual detection subsystem achieved its primary objectives. 
The YOLOv8n model trained on the drone detection dataset of 
Pawełczyk and Wojtyra achieved an mAP@0.50 of 0.835 and an F1 
score of 0.81 at a confidence threshold of 0.571, outperforming 
the MobileNet V1 baseline from the original dataset paper. 
Deployed on a Raspberry Pi~5 with a Hailo AI HAT+ accelerator, 
the model achieved approximately 108 frames per second inference, 
well within the requirements for real-time tracking, at a CPU 
usage of only 29.7\% leaving substantial headroom for system 
integration. Real-time coordinate publishing over MQTT and live 
video streaming to the web-based GUI were both successfully 
demonstrated. The full closed-loop camera tracking, where motor 
movement is driven continuously by YOLO detections following the 
initial acoustic orientation, was not completed within the project 
timeframe due to issues with the motor control hardware, which 
is identified as the primary target for future development.

The complete system was realised at a total hardware cost of 
approximately NOK~8,503, representing a reduction of over an order 
of magnitude compared to entry-level commercial drone detection 
solutions, while demonstrating real-time multimodal detection 
capability. The coarse-to-fine architecture, in which acoustic 
DOA estimation narrows the search space for the visual detector, 
proved sound in principle and was partially validated through 
separate testing of each subsystem.

Future work should prioritise completing the closed-loop motor 
tracking, conducting outdoor detection performance evaluation 
at measured distances using consumer quadcopters in 
representative civilian environments, and improving the 
decision-tree classifier's generalisation to real microphone 
array data. The addition of proper quantization calibration 
using representative drone images during the Hailo HEF 
conversion process would also likely improve visual detection 
accuracy on the deployed hardware.